% =========================================================
%  PRESENTATION TEMPLATE — Theme 3B, Mathematical Formalisation with Lean
%
%  Using this is optional. If you already have a way of making
%  slides you are happy with, use that. This is only here to save you
%  the setup time.
%
%  Compile with xelatex, not pdflatex. The Lean symbols need it.
%  There is no LaTeX in the Codespace, so the easy route is Overleaf:
%  sign in at overleaf.com/edu/bristol, upload this file, then set
%  Menu -> Compiler -> XeLaTeX.
%
%  Uses the `metropolis' beamer theme, which Overleaf, TeX Live and
%  MacTeX all have.
%
%  There is no running order here. How you build the talk is up to
%  you. The slides at the end show what you can use; delete them.
%
%  Change anything you like: colours, theme, fonts, the boxes.
%
%  A frame with Lean code in it must say [fragile].
% =========================================================
\documentclass[11pt,aspectratio=169]{beamer}
\usetheme[progressbar=frametitle]{metropolis}
\setbeamertemplate{navigation symbols}{}

% ---- Colours. Change these and the whole talk changes. ----
\usepackage{xcolor}
\definecolor{talkmain}{HTML}{2F6FB3}
\definecolor{talkaccent}{HTML}{F2A900}
\definecolor{talktext}{HTML}{22303A}
\setbeamercolor{frametitle}{bg=talkmain, fg=white}
\setbeamercolor{progress bar}{fg=talkaccent, bg=talkaccent!25}
\setbeamercolor{title separator}{fg=talkaccent}
\setbeamercolor{normal text}{fg=talktext}
\setbeamercolor{alerted text}{fg=talkmain}

\usepackage{fontspec}
\setmonofont{DejaVu Sans Mono}[Scale=0.8]
\usepackage{amsmath,amssymb}
\usepackage{booktabs}
\usepackage{listings}
\usepackage[most]{tcolorbox}

% ---- Lean code ----
\definecolor{keywordcolor}{rgb}{0.4,0.1,0.5}
\definecolor{commentcolor}{rgb}{0.3,0.5,0.3}
\lstdefinelanguage{lean}{
  keywords={theorem,lemma,by,induction,intro,exact,apply,rw,rfl,simp,ring,
    linarith,omega,decide,sorry,open,import,namespace,end},
  sensitive=true, comment=[l]{--}, morecomment=[s]{/-}{-/},
  keywordstyle=\color{keywordcolor}\bfseries,
  commentstyle=\color{commentcolor}\itshape,
}
\lstset{language=lean, basicstyle=\ttfamily\footnotesize, columns=fullflexible,
  keepspaces=true, breaklines=true}

\newcommand{\lean}[1]{\texttt{#1}}

% ---- Three boxes, ready to use. Examples at the end of the file. ----
% \begin{keybox}{Title}  ... \end{keybox}   — the point of the slide
\newtcolorbox{keybox}[1]{enhanced, colback=talkmain!6, colframe=talkmain,
  boxrule=0.8pt, arc=1mm, left=2mm, right=2mm, top=1.2mm, bottom=1.2mm,
  title={#1}, coltitle=white, colbacktitle=talkmain, fonttitle=\bfseries\small}
% \begin{notebox} ... \end{notebox}         — an aside, no title
\newtcolorbox{notebox}{enhanced, colback=talkaccent!12, colframe=talkaccent!70!black,
  boxrule=0.6pt, arc=1mm, left=2mm, right=2mm, top=1.2mm, bottom=1.2mm}
% \begin{plainbox} ... \end{plainbox}       — a quiet frame around anything
\newtcolorbox{plainbox}{enhanced, colback=white, colframe=black!25,
  boxrule=0.5pt, arc=1mm, left=2mm, right=2mm, top=1.2mm, bottom=1.2mm}

\title{<Title of your result>}
\subtitle{<one line saying what you formalised>}
\author{<Name> \and <Name> \and <Name>}
\institute{Perspectives in Mathematics --- Theme 3B}
\date{}

\begin{document}

\begin{frame}
  \titlepage
\end{frame}

% ---------------------------------------------------------
%  Your slides go here. One slide is
%
%      \begin{frame}{Title}  ...your content...  \end{frame}
%
%  Leave the braces empty for a slide with no title. A slide with
%  Lean code on it must say [fragile]:
%
%      \begin{frame}[fragile]{Title}
%
%  Add [fragile] when you paste the code in, not before --- an empty
%  [fragile] frame gives a strange LaTeX error.
% ---------------------------------------------------------

% =========================================================
%  The slides below are just a show of what you can use.
%  Copy what is useful and DELETE THEM.
% =========================================================

\begin{frame}{What you can use (delete these slides)}
  Three boxes. The colours are only a suggestion:
  \begin{keybox}{Theorem}
    For every $n \in \mathbb{N}$, \quad
    $\displaystyle\sum_{k=0}^{n-1}(2k+1) = n^2$.
  \end{keybox}
  \begin{notebox}
    In $\mathbb{N}$, subtraction truncates: $2 \cdot 0 - 1 = 0$.
  \end{notebox}
  \begin{plainbox}
    \centering A quiet frame around a picture, or anything else.
  \end{plainbox}
\end{frame}

% Note the [fragile] below: any slide holding Lean code needs it.
\begin{frame}[fragile]{Lean code}
\begin{lstlisting}
theorem my_result (n : ℕ) : ∑ k ∈ range n, (2 * k + 1) = n ^ 2 := by
  induction n with
  | zero => simp
  | succ d hd => sorry
\end{lstlisting}
  A Lean name inside a sentence: \lean{Finset.sum\_range\_succ}.
\end{frame}

\begin{frame}{Two columns, and a table}
  \begin{columns}[t]
    \column{0.5\textwidth}
      \textbf{On paper.} The induction reduces to one line of algebra.
    \column{0.5\textwidth}
      \textbf{In Lean.} The same step, with every type made explicit.
  \end{columns}
  \vspace{1em}
  \begin{center}
  \small
  \begin{tabular}{@{}lll@{}}
  \toprule
  \textbf{Declaration} & \textbf{Lead} & \textbf{Built or reused?} \\
  \midrule
  \lean{step}      & <name> & built (\lean{ring}) \\
  \lean{main\_thm} & <name> & built on \lean{sum\_range\_succ} \\
  \bottomrule
  \end{tabular}
  \end{center}
\end{frame}

\begin{frame}{}
  \begin{center}
    {\Large <one big line: the sentence you want them to remember>}
  \end{center}
\end{frame}

\end{document}

